\documentclass[reqno]{amsart} \input{common.tex}

\begin{document}

\title{Siegel transforms and counting rational approximations on flag varieties}

\begin{abstract}
  Talk by Ren\'e Pfitscher (Paris 13).
\end{abstract}

Dirichlet: infinitely many rational approximations with
\begin{equation*}
  \left| x - \frac{p}{q} \right| < \frac{1}{q^2}.
\end{equation*}
Khintchine: let $\psi : \mathbb{N} \rightarrow \mathbb{R}_+$ be non-increasing.  Then for almost all $x \in \mathbb{R}$, the inequality
\begin{equation*}
  \left| x - \frac{p}{q} \right| < \psi(q)
\end{equation*}
has infinitely (resp.\ finitely) many solutions according as the sum $\sum_{q \geq 1} \psi(q) q$ is infinite (resp.\ finite).

Schmidt, 1960: Set
\begin{equation*}
  N_\psi(x, T) := \# \left\{(p, q) \in \mathbb{Z} \times \mathbb{N} :(\ast) \text{ holds, } q < T \right\}.
\end{equation*}
Then for almost all $x \in \mathbb{R}$, we have
\begin{equation*}
  N_\psi(x, T) \sim 2 \sum_{q = 1}^T \psi(q) q \text{ as } T \rightarrow \infty.
\end{equation*}

Now let's say a bit about \emph{flag varieties}.  More generally, let $G$ be a semisimple $\mathbb{Q}$-group, and $P$ a parabolic $\mathbb{Q}$-subgroup.  We set $X := G / P$.  We choose a maximal compact subgroup $K$.  We also need a height function $H_\chi : X(\mathbb{Q}) \rightarrow \mathbb{R}_+$.  We also need a $K$-invariant Riemannian distance function $d(\cdot, \cdot)$.  We denote by $\sigma$ the corresponding $K$-invariance Riemannian probability measure, given by the normalized volume form.

For example, we can take
\begin{equation*}
  G = \SL_2(\mathbb{R}) > P =
  \begin{pmatrix}
    \ast    & \ast \\
    0         & \ast \\
  \end{pmatrix}, \quad
  K = \SO(2),
\end{equation*}
so that $X = \mathbb{P}^1$, and
\begin{equation*}
  H : \mathbb{P}^1(\mathbb{Q}) \rightarrow \mathbb{R}_+
\end{equation*}
given as follows: given $v \in \mathbb{P}^1(\mathbb{Q})$, we can represent it by a primitive lattice point $\vec{v} \in \mathbb{Z}^2_{\mathrm{prim}}$, and we then define the height to be the norm of that vector:
\begin{equation*}
  H(v) := \lVert \vec{v} \rVert.
\end{equation*}

We can picture this in terms of shortest integer vectors along rational lines.

Set
\begin{equation*}
  N_\tau(x, T) = \# \left\{ v \in X(\mathbb{Q}) : d(x, v < H_\chi(v)^{- \tau}), \, H_\chi(v) < T \right\},
\end{equation*}
\begin{equation*}
  \mathcal{E}(x, T) = \left\{ \vec{v} \in \mathbb{R}^2 - \{0\} : \lVert \vec{v} \rVert < T,
    \,
    d(x,[\vec{v}]) < \lVert \vec{v} \rVert^{- \tau}
  \right\}.
\end{equation*}
In the example of interest,
\begin{equation*}
  N_\tau(x, T) = \#(\mathbb{Z}_{\mathrm{prim}}^2 \cap \mathcal{E}(x, T))
  =
  \sum_{v \in \mathbb{Z}^2_{\mathrm{prim}}} 1_{\mathcal{E}(x, T)}(v).
\end{equation*}

We form the dyadic region
\begin{equation*}
  \mathcal{F}(x, T) := \mathcal{E}(x, T) - \mathcal{E}(x, T/2).
\end{equation*}
We can apply a flow, such as $a_y := \diag(y^{-1}, y)$, to this region.  We choose $y_T$ so that the flowed region, $\mathcal{B}(T)$, is nice and round.  Then
\begin{align*}
  \#(\mathbb{Z}^2_{\mathrm{prim}} \cap \mathcal{F}(x, T)) &= \# \left( k_X \mathbb{Z}^2_{\mathrm{prim}} \cap \mathcal{F}(T) \right)
  \\ &=
  \# \left( a_{y_T} k_X \mathbb{Z}^2_{\mathrm{prim}} \cap \mathcal{B}(T) \right) \\
  &\approx \vol(\mathcal{B}(T)).
\end{align*}
Here $\mathcal{B}(T)$ ``looks like a square'', while $k_X$ is typically not ``too distorted''.

\begin{lemma}
  There exists $\eps > 0$ such that for all large $T$, and all $g \in G$ such that $\operatorname{height}(g) \ll T^\eps$, where
  \begin{equation*}
    \operatorname{height}(g) := \left( \min_{v \in \mathbb{Z}^2 - 0} \lVert g v \rVert \right)^{-1},
  \end{equation*}
  we have
  \begin{equation*}
    \# \left( g \mathbb{Z}^2_{\mathrm{prim}} \cap \mathcal{B}(T) \right)
    = \vol(\mathcal{B}(T)) \left( 1 + \O(T^{- \eps}) \right).    
  \end{equation*}
\end{lemma}
\begin{proof}[Proof sketch]
  This is a counting argument due to (or at least inspired by) Eskin--McMullen.

  Let's first notice that what we want to estimate is, as we have said before, the primitive Siegel transform
  \begin{equation*}
    F_T(g) := \#(g \mathbb{Z}^2_{\mathrm{prim}} \cap \mathcal{B}(T)) = \sum_{v \in \mathbb{Z}^2_{\mathrm{prim}}} 1_{\mathcal{B}(T)}(g v)
    = \sum_{\gamma \in \Gamma / \Gamma_\infty} 1_{\mathcal{B}(T)}(g \gamma e_1).
  \end{equation*}
  Here $\Gamma_\infty = \Gamma \cap U$, where $U := \left(
    \begin{smallmatrix}
      1&\ast\\
      0       &1 \\
    \end{smallmatrix}
  \right)$.  Now we unfold:
  \begin{align*}
    \langle F_T, \phi \rangle &= \int_{G / \Gamma} F_T(g) \phi(g) \, d \mu(g) =
    \int_{G / \Gamma} \sum_{\gamma \in \Gamma / \Gamma_\infty} 1_{\mathcal{B}(T)}
                 (g \gamma e_1) \phi(g) \, d \mu  \\
               &= \int_{G / U} \int_{U / \Gamma_\infty}
                 1_{\mathcal{B}(T)}(g u e_1) \phi(g u)
                 \, d u \, d m \\
               &=
                 \int_{\mathbb{R}^2} 1_{\mathcal{B}(T)}(k a_y e_1)
                 \int_{U / \Gamma_\infty}
                 \phi(k a_y u) \, d u \, d m.
  \end{align*}
  Now we use the equidistribution result of Sarnak (1981) to see that there exists $\alpha > 0$ so that for all $f \in C_c^\infty(G / \Gamma)$, we have
  \begin{equation*}
    \left| \int_{U / \Gamma_\infty} \phi(k a_y u) \, d u - \int_{G / \Gamma} \phi \, d \mu \right|
    \ll y^{- \alpha} S(\phi).    
  \end{equation*}
  This yields the required asymptotic.  The well-roundedness allows us to go from an average counting result to an individual counting result.  Main term: $f = 1_{\mathcal{B}(T)}$,
  \begin{equation*}
    \vol(\mathcal{B}(T)) = \int_{\mathbb{R}^2} f \, d m
    = \int_{G / \Gamma} \hat{f} \, d \mu.
  \end{equation*}
  Variance:
  \begin{equation*}
    \int_{G / \Gamma} \left| \hat{f} - \vol(\mathcal{B}(T)) \right|^2 \, d \mu
    =
    \int_{G / \Gamma} \left| \hat{f} \right|^2 \, d \mu
    -
    \vol(\mathcal{B}(T))^2.    
  \end{equation*}
  The integral on the right hand side may be written
  \begin{equation*}
    \langle \hat{f}, \hat{f} \rangle =
    \int_{\mathbb{R}^2} 1_{\mathcal{B}(T)}(k a_y e_1)
    \int_{U / \Gamma_\infty}
    E(k a_y u, f) \, d u \, d m.
  \end{equation*}
  This last integral is a constant term of an Eisenstein series.  It describes second moments of Siegel transforms.  One obtains in this way a bound on the variance.  One can also argue and get an asymptotic estimate.
\end{proof}

Okay, here's our result:
\begin{theorem}[Pfitscher 2024]
  Suppose that $P$ is a \emph{maximal} parabolic subgroup of any semisimple algebraic $\mathbb{Q}$-group $G$.  Let $\tau \geq 0$ be arbitrary.  Write $d := \dim X$, let $\beta_\chi$ denote the Diophantine exponent, and define
  \begin{equation*}
    \Psi(T) :=
    \begin{cases}
      T^{(\beta_\chi - \tau) d} & \text{ if } \tau < \beta_\chi, \\
      \log T& \text{ if } \tau = \beta_\chi.
    \end{cases}
  \end{equation*}
  Then there exists $c > 0$ and $\eps > 0$ such that for almost every $x \in X$, the counting function is given by
  \begin{equation*}
    \mathcal{N}_\tau(x, T) = c \Psi(T) \left( 1 + \O_X(\Psi(T)^{- \eps}) \right).
  \end{equation*}
\end{theorem}

\begin{remark}
  For $\mathbb{R}^n$, you can prove that the second moment of the Siegel transform is given as follows:
  \begin{equation*}
    \int_{G / \Gamma} \lvert \hat{f} \rvert^2 \, d \mu -
    \left( \int_{\mathbb{R}^n} f \, d m \right)^2 = \int_{\mathbb{R}} (\lvert f \rvert^2 +
    f \cdot f^-) \, d m \ll \vol(\mathcal{B}(T)),
  \end{equation*}
  where
  \begin{equation*}
    f^-(v) = f(- v).
  \end{equation*}
  This estimate suffices to get the conclusion of the theorem (for an arbitrary decreasing $\Psi$).  Kelmer--Yu gave an analogue for the sphere.  Our proof used dynamics.
\end{remark}

\bibliography{refs}{} \bibliographystyle{plain}
\end{document}
